%% Dokument 168 – Kapitel (Deutsche Version)
%% Das Periodensystem als ξ-Geometrie:
%% Radiale Massenleiter und angulare Schalenstruktur aus einem gemeinsamen Ursprung
%% Johann Pascher, März 2026

\chapter{Das Periodensystem als \texorpdfstring{$\xi$}{xi}-Geometrie}
\label{chap:periodensystem-xi}

\begin{abstract}
Das Periodensystem der Elemente und die Teilchenmassen-Leiter der T0-Theorie
teilen dieselbe geometrische Quantenzahl~$n$, jedoch in verschiedenen Rollen.
Die Teilchenmassen folgen einer \emph{radialen} exponentiellen Skalierung
$m_n \sim \xi^{-n} \cdot m_P$, während die Schalenstruktur des Periodensystems
einer \emph{angularen} quadratischen Entartung $g_n = 2n^2$ gehorcht.
Beide entstehen aus der dreidimensionalen Gittergeometrie der T0-Theorie:
$\xi$ ist das Packungsdefizit des 3D-Gitters (radiale Richtung),
$2n^2$ ist der Entartungsgrad der $n$-ten Kugelschale im selben Gitter
(angulare Richtung).
Die Rydberg-Energie liegt exakt auf Stufe~7 der $\xi$-Potenzleiter:
$R_\infty = \frac{m_e \alpha^2}{2} \approx \xi^{6{,}956} \cdot E_P$.
Das Dualitätsprinzip $E_n \cdot g_n = 2R_\infty = \text{const}$
verbindet die $1/n^2$-Energieleiter mit der $n^2$-Füllleiter
und ist eine geometrische Äquipartition.
Das Periodensystem ist der angulare Querschnitt der $\xi$-Geometrie
auf Stufe $\xi^{\sim 7}$ — analog dazu sind die Leptonmassen der radiale
Querschnitt auf den Stufen $\xi^{4{,}9}$ bis $\xi^{5{,}8}$.
\end{abstract}

%──────────────────────────────────────────────────────────────────────────────
\section{Einleitung und Fragestellung}
\label{sec:ps-einleitung}

Die T0-Theorie (Zeit-Masse-Dualität / FFGFT) leitet alle physikalischen
Konstanten aus dem geometrischen Parameter
\begin{equation}
  \xi = \frac{4}{30\,000} = 1{,}3\overline{3} \times 10^{-4}
\end{equation}
ab \cite{dok009,dok041,dok056}.
Die Teilchenmassen des Standardmodells organisieren sich auf einer
exponentiellen Leiter
\begin{equation}
  m_n \;\sim\; \xi^{-n} \cdot m_P
  \label{eq:massenleiter}
\end{equation}
wobei $n$ die T0-Quantenzahl und $m_P$ die Planck-Masse ist.

Das Periodensystem der Elemente hingegen zeigt ein ganz anderes Muster:
Die Schalenkapazitäten folgen
\begin{equation}
  g_n \;=\; 2n^2 \;=\; 2,\; 8,\; 18,\; 32,\; 50,\;\ldots
  \label{eq:schalen}
\end{equation}
und die Bindungsenergien skalieren mit $1/n^2$.
Auf den ersten Blick scheinen Teilchenmassen und Periodensystem unverbunden.

\textbf{Geminis Beobachtung} (2026): Das Periodensystem liegt „eine Ebene
höher" als die Teilchenmassen — beide haben Quantenzahlen, aber welcher
gemeinsame geometrische Ursprung verbindet die Muster?

Die Antwort dieses Dokuments: \emph{Die T0-Geometrie hat zwei Achsen —
eine radiale ($\xi^n$-Skalierung) und eine angulare ($2n^2$-Entartung),
die beide aus der dreidimensionalen Gitterstruktur mit Packungsdefizit~$\xi$
folgen.}

%──────────────────────────────────────────────────────────────────────────────
\section{Die \texorpdfstring{$\xi$}{xi}-Potenzleiter und ihre Stufen}
\label{sec:xi-leiter}

Die bisherigen Dokumente der T0-Reihe haben die $\xi$-Potenzleiter
schrittweise aufgebaut. Tabelle~\ref{tab:xi-leiter} zeigt den Stand
inklusive der neuen atomaren Stufe.

\begin{table}[htbp]
\centering
\caption{Vollständige $\xi$-Potenzleiter (Stand: März 2026)}
\label{tab:xi-leiter}
\begin{tabular}{clll}
\hline
Stufe & $\xi$-Potenz & Physikalischer Inhalt & Dokument \\
\hline
$0$       & $\xi^0 = 1$                & Planck-Skala, $L_0 = \xi \cdot \ell_P$ & Dok.\,009 \\
$1/2$     & $\xi^{1/2} \approx 10^{-2}$ & $\alpha_W$, $m_e/m_\mu$, PPLN $\Delta n$ & Dok.\,041, 167 \\
$1$       & $\xi^1 \approx 10^{-4}$    & $\alpha_\text{EM} \sim \xi \cdot E_0^2$  & Dok.\,011 \\
$4$       & $\xi^4 \approx 10^{-16}$   & CMB-Temperatur $T_\text{CMB}$            & Dok.\,166 \\
$23/4$    & $\xi^{5{,}775}$            & Elektronmasse $m_e/m_P$                  & Dok.\,056 \\
$\mathbf{\sim 7}$ & $\boldsymbol{\xi^{6{,}956}}$ & \textbf{Rydberg-Energie $R_\infty$, Periodensystem} & \textbf{dieses Dok.} \\
$10$      & $\xi^{10} \approx 10^{-43}$& Hubble-Konstante $H_0$                   & Dok.\,165 \\
\hline
\end{tabular}
\end{table}

Die Rydberg-Energie liegt zwischen den Stufen~6 und~7:
\begin{equation}
  \frac{R_\infty}{E_P}
  = \frac{m_e}{m_P} \cdot \frac{\alpha^2}{2}
  = \xi^{5{,}775} \cdot \xi^{1{,}181}
  = \xi^{6{,}956}
  \approx \xi^7
\end{equation}
Der Exponent setzt sich aus zwei unabhängigen T0-Größen zusammen:
der Elektronmasse (Stufe $5{,}775$, Dok.\,056) und der Feinstrukturkonstante
(Stufe $1{,}181$, Dok.\,011).

%──────────────────────────────────────────────────────────────────────────────
\section{Zwei Achsen der \texorpdfstring{$\xi$}{xi}-Geometrie}
\label{sec:zwei-achsen}

\subsection{Radiale Achse: Exponentielle Massenleiter}
\label{subsec:radial}

Die Teilchenmassen organisieren sich auf der \emph{radialen} Achse des
T0-Gitters. Die Grundrelation lautet \cite{dok009}:
\begin{equation}
  m_n \;\sim\; \xi^{-n} \cdot m_P, \qquad n = 1, 2, 3, \ldots
\end{equation}
Jeder Schritt auf dieser Leiter multipliziert die Masse mit
$\xi^{-1} = 7500$.
Die Leptonen belegen die Stufen (Dok.\,056, 122):
\begin{align}
  m_e:  &\quad \xi^{5{,}775} \cdot m_P \\
  m_\mu:&\quad \xi^{5{,}177} \cdot m_P \\
  m_\tau:&\quad \xi^{4{,}861} \cdot m_P
\end{align}
Das Massenverhältnis $m_e/m_\mu \propto \xi^{1/2}$ (Dok.\,041) liegt
exakt auf dem elektroschwachen Exponenten — auf der Halbstufe der
$\xi$-Leiter.

\subsection{Angulare Achse: Quadratische Entartungsleiter}
\label{subsec:angular}

Die \emph{angulare} Achse desselben Gitters bestimmt die Entartungsgrade
der Kugelschalen. Im 3D-Raum gilt für die $n$-te Hauptschale:
\begin{equation}
  g_n \;=\; 2 \sum_{l=0}^{n-1} (2l+1) \;=\; 2n^2
  \label{eq:gn}
\end{equation}
Der Faktor 2 kommt vom Spin $m_s = \pm\frac{1}{2}$, die Summe ist reine
Drehimpuls-Algebra im 3D-Raum. Tabelle~\ref{tab:schalen} zeigt die ersten
vier Schalen.

\begin{table}[htbp]
\centering
\caption{Schalenstruktur und Periodensystem}
\label{tab:schalen}
\begin{tabular}{ccccc}
\hline
$n$ & $g_n = 2n^2$ & Periodenlänge & Edelgas~$Z$ & Periodennummer \\
\hline
1 &  2 &  2 &  He (2)  & 1 \\
2 &  8 &  8 &  Ne (10) & 2, 3 \\
3 & 18 & 18 &  Ar (18) & 4, 5 \\
4 & 32 & 32 &  Kr (36) & 6, 7 \\
\hline
\end{tabular}
\end{table}

\subsection{Die Verbindung durch die Quantenzahl \texorpdfstring{$n$}{n}}
\label{subsec:verbindung-n}

Dieselbe Quantenzahl~$n$ erscheint in beiden Systemen, aber mit
verschiedener funktionaler Rolle:

\begin{align}
  \text{Radial (Massen):}  &\quad m_n \;\sim\; \xi^{-n} \cdot m_P
    \quad\text{(exponentiell)} \\
  \text{Angular (Schalen):}&\quad g_n \;=\; 2n^2
    \quad\text{(quadratisch)}
\end{align}

Das ist kein Zufall. In der T0-Gittergeometrie entspricht~$n$ dem
\emph{diskreten Radius im Gitter}:
\begin{itemize}
  \item Radial: $n$ Schritte aus dem Gitter heraus $\Rightarrow$
    Masse skaliert mit $\xi^{-n}$ (geometrische Folge)
  \item Angular: die Gitterpunkte auf der Sphäre mit Radius~$n$ haben
    Entartungsgrad $2n^2$ (kombinatorische Folge)
\end{itemize}

%──────────────────────────────────────────────────────────────────────────────
\section{Das Dualitätsprinzip}
\label{sec:dualitaet}

\subsection{Energieleiter und Füllleiter sind dual}
\label{subsec:dual}

Auf der atomaren Skala (Wasserstoff-ähnlich) gilt:
\begin{align}
  \text{Energieleiter:} &\quad E_n \;=\; -\frac{R_\infty}{n^2}
    \quad\text{(fällt wie } 1/n^2\text{)} \\
  \text{Füllleiter:}    &\quad g_n \;=\; 2n^2
    \quad\text{(wächst wie } n^2\text{)}
\end{align}

Das Produkt ist exakt konstant:
\begin{equation}
  \boxed{E_n \cdot g_n \;=\; \frac{R_\infty}{n^2} \cdot 2n^2 \;=\; 2R_\infty
  \;=\; \text{const}}
  \label{eq:dualitaet}
\end{equation}

Dies ist geometrische \emph{Äquipartition}: jede Schale trägt dieselbe
Gesamtenergie zur Rydberg-Skala bei, unabhängig von~$n$.
Die $1/n^2$-Energieleiter und die $n^2$-Füllleiter sind \emph{invers dual}.

\subsection{Vergleich mit der Massenleiter}

Tabelle~\ref{tab:vergleich} fasst die drei Leitern zusammen.

\begin{table}[htbp]
\centering
\caption{Drei Leitern — ein geometrischer Ursprung}
\label{tab:vergleich}
\begin{tabular}{lccc}
\hline
System & Leiter & Funktion von $n$ & Typ \\
\hline
Teilchenmassen & $m_n \sim \xi^{-n} \cdot m_P$      & exponentiell & radial \\
Schalenenergie & $E_n = R_\infty / n^2$              & harmonisch   & angular \\
Schalenfüllung & $g_n = 2n^2$                        & quadratisch  & angular \\
\hline
\multicolumn{4}{l}{Dualität: $E_n \cdot g_n = 2R_\infty = \text{const}$} \\
\hline
\end{tabular}
\end{table}

%──────────────────────────────────────────────────────────────────────────────
\section{Die Rolle der fraktalen Dimension}
\label{sec:fraktal}

In T0 gilt (Dok.\,009, 133):
\begin{equation}
  D_f \;=\; 3 - \xi \;=\; 2{,}999867
\end{equation}
Die Dimension des Raumes ist fast, aber nicht exakt~3.

Im Grenzfall $D_f \to 3$ ($\xi \to 0$):
\begin{itemize}
  \item $g_n = 2n^2$ ist \emph{exakt} (reine Kugelsymmetrie im $\mathbb{Z}^3$-Gitter)
  \item die Schalenstruktur des Periodensystems ist exakt
\end{itemize}

Für $D_f = 3 - \xi < 3$ gibt es eine minimale Korrektur:
\begin{equation}
  \delta g_n \;\sim\; \xi \cdot n^2
\end{equation}
Für $n = 2$: $\delta g_2 \approx \xi \cdot 4 \approx 5 \times 10^{-4}$
(völlig vernachlässigbar).

Das Periodensystem ist also exakt, weil $D_f \approx 3$ auf 4 Dezimalstellen.
Die winzige Abweichung $\xi \approx 1{,}3 \times 10^{-4}$ macht die Raumgeometrie
„fast-ganzzahlig" — und genau diese Abweichung generiert die gesamte
Physik der Teilchenmassen und Fundamentalkonstanten.

%──────────────────────────────────────────────────────────────────────────────
\section{Numerische Überprüfung}
\label{sec:numerik}

\subsection{Rydberg-Energie auf der \texorpdfstring{$\xi$}{xi}-Leiter}

\begin{align}
  R_\infty &= \frac{m_e \alpha^2}{2} = 13{,}5984\;\text{eV} \\
  \frac{R_\infty}{E_P} &= \frac{m_e}{m_P} \cdot \frac{\alpha^2}{2}
    = \xi^{5{,}775} \cdot \xi^{1{,}181}
    = \xi^{6{,}956}
\end{align}

\begin{table}[htbp]
\centering
\caption{Rydberg-Energie: T0-Ableitung vs. Messung}
\label{tab:rydberg}
\begin{tabular}{lcc}
\hline
Größe & Wert & $\xi$-Exponent \\
\hline
$m_e / m_P$       & $4{,}185 \times 10^{-23}$ & $5{,}775$ \\
$\alpha^2 / 2$    & $2{,}663 \times 10^{-5}$  & $1{,}181$ \\
$R_\infty / E_P$  & $1{,}114 \times 10^{-27}$ & $6{,}956$ \\
$R_\infty$ (T0)   & $\approx 9{,}1\;\text{eV}$ (Näherung $\xi^7$) & $7{,}000$ \\
$R_\infty$ (gemessen) & $13{,}598\;\text{eV}$ & $6{,}956$ \\
\hline
\end{tabular}
\end{table}

Der exakte Exponent $6{,}956 \approx 7$ zeigt, dass die Rydberg-Energie
auf Stufe~7 der $\xi$-Leiter liegt. Die Abweichung von der ganzen Zahl~7
ergibt sich aus dem exakten T0-Wert $m_e/m_P = \xi^{5{,}775}$ (nicht $\xi^6$).

\subsection{Dualitätsrelation: numerische Verifikation}

\begin{align}
  E_1 \cdot g_1 &= 13{,}60\;\text{eV} \cdot 2  = 27{,}20\;\text{eV} = 2R_\infty \\
  E_2 \cdot g_2 &= \phantom{0}3{,}40\;\text{eV} \cdot 8  = 27{,}20\;\text{eV} = 2R_\infty \\
  E_3 \cdot g_3 &= \phantom{0}1{,}51\;\text{eV} \cdot 18 = 27{,}18\;\text{eV} \approx 2R_\infty \\
  E_4 \cdot g_4 &= \phantom{0}0{,}85\;\text{eV} \cdot 32 = 27{,}20\;\text{eV} = 2R_\infty
\end{align}

Die Dualitätsrelation~\eqref{eq:dualitaet} gilt exakt für alle~$n$.

%──────────────────────────────────────────────────────────────────────────────
\section{Das Periodensystem in der T0-Hierarchie}
\label{sec:hierarchie}

\begin{figure}[htbp]
\centering
\begin{tabular}{|l|l|l|}
\hline
\textbf{$\xi$-Stufe} & \textbf{Physik} & \textbf{Achse} \\
\hline
$\xi^0 = 1$      & Planck-Skala ($L_0 = \xi \cdot \ell_P$) & — \\
$\xi^{1/2}$      & Elektroschwach: $\alpha_W$, $m_e/m_\mu$  & radial \\
$\xi^1$          & $\alpha_\text{EM}$, EM-Kopplung           & radial \\
$\xi^4$          & CMB-Temperatur $T_\text{CMB}$             & radial \\
$\xi^{5{,}8}$    & Teilchenmassen (Leptonen, Quarks)         & radial \\
$\xi^{6{,}956}$  & \textbf{Rydberg $R_\infty$, Periodensystem} & \textbf{beide} \\
$\xi^{10}$       & Hubble-Konstante $H_0$                    & radial \\
\hline
\end{tabular}
\caption{Position des Periodensystems in der vollständigen $\xi$-Hierarchie}
\label{fig:hierarchie}
\end{figure}

\vspace{0.5em}

Die Besonderheit der atomaren Stufe $\xi^{\sim 7}$: hier wirken
\emph{beide} Achsen. Die radiale Achse bestimmt die Energieskala $R_\infty$,
die angulare Achse bestimmt das Muster $2n^2$. Das Periodensystem ist
das sichtbare Resultat beider Achsen gemeinsam.

%──────────────────────────────────────────────────────────────────────────────
\section{Vorhersagen und offene Fragen}
\label{sec:vorhersagen}

\subsection{Angulare Korrekturen durch \texorpdfstring{$D_f < 3$}{Df<3}}

Die fraktale Dimension $D_f = 3 - \xi$ führt zu minimalen Korrekturen
der exakten $2n^2$-Regel:
\begin{equation}
  g_n^\text{(T0)} \;=\; 2n^2 \cdot (1 - \xi \cdot f(n))
\end{equation}
wobei $f(n)$ eine dimensionslose Funktion von~$n$ ist (noch abzuleiten).
Für $n \leq 4$ ist die Korrektur $< 10^{-3}$ und nicht messbar.

\subsection{Bindungsenergien molekularer Systeme}

Die chemischen Bindungsenergien liegen bei $1$–$10\;\text{eV}$ und damit
ebenfalls auf Stufe~$\xi^7$. Es wäre zu prüfen, ob spezifische Bindungsenergien
auf rationalen Vielfachen von $\xi^{7/2} \cdot R_\infty$ liegen.

\subsection{Schwere Elemente und relativistische Korrekturen}

Für $Z \gg 1$ werden relativistische Effekte wichtig. In T0 hängen diese
mit der $\xi^{1/2}$-Elektroschwachstufe zusammen. Eine systematische
Analyse der Ionisierungsenergien schwerer Elemente ($Z > 54$) könnte
$\xi$-Treffer auf der halben Stufe $\xi^{7-1/2} = \xi^{13/2}$ zeigen.

%──────────────────────────────────────────────────────────────────────────────
\section{Zusammenfassung}
\label{sec:ps-zusammenfassung}

\begin{enumerate}
  \item \textbf{Energieskala:} $R_\infty / E_P = \xi^{6{,}956} \approx \xi^7$.
    Die Rydberg-Energie — und damit die gesamte Atomphysik — liegt auf
    Stufe~7 der $\xi$-Potenzleiter.

  \item \textbf{Radiale Achse:} Teilchenmassen folgen $m_n \sim \xi^{-n} \cdot m_P$.
    Der Exponent~$n$ ist der diskrete Gitterradius in der radialen $\xi$-Richtung.

  \item \textbf{Angulare Achse:} Schalenkapazitäten folgen $g_n = 2n^2$.
    Der Index~$n$ ist derselbe diskrete Gitterradius, jetzt als Schalen-Index
    in der angularen 3D-Sphären-Richtung.

  \item \textbf{Dualitätsprinzip:} $E_n \cdot g_n = 2R_\infty = \text{const}$.
    Die $1/n^2$-Energieleiter und die $n^2$-Füllleiter sind invers dual.
    Jede Schale trägt dieselbe Gesamtenergie $2R_\infty$ zur $\xi^7$-Skala bei.

  \item \textbf{Fraktale Dimension:} $D_f = 3 - \xi \approx 3$ macht $2n^2$ exakt
    (auf $\sim 10^{-4}$). Das Periodensystem ist exakt, weil der Raum
    fast-ganzzahlig-dimensional ist.

  \item \textbf{Gemeinsamer Ursprung:} Teilchenmassen und Periodensystem
    entstehen aus derselben 3D-Gittergeometrie mit Packungsdefizit~$\xi$.
    Das Periodensystem ist der \emph{angulare Querschnitt} der $\xi$-Geometrie
    auf Stufe $\xi^{\sim 7}$; die Leptonmassen sind der \emph{radiale Querschnitt}
    auf Stufen $\xi^{4{,}9}$–$\xi^{5{,}8}$.
\end{enumerate}

\begin{thebibliography}{99}
\bibitem{dok009} Pascher, J.: \emph{Ursprung des $\xi$-Parameters (Dok.\,009)},
  T0-Theorie-Reihe, 2025.
\bibitem{dok011} Pascher, J.: \emph{Feinstrukturkonstante (Dok.\,011)},
  T0-Theorie-Reihe, 2025.
\bibitem{dok041} Pascher, J.: \emph{Parameterherleitung (Dok.\,041)},
  T0-Theorie-Reihe, 2025.
\bibitem{dok056} Pascher, J.: \emph{Universale Ableitung der Leptonmassen (Dok.\,056)},
  T0-Theorie-Reihe, 2025.
\bibitem{dok122} Pascher, J.: \emph{$m_e/m_\mu$-Verhältnis absolut (Dok.\,122)},
  T0-Theorie-Reihe, 2025.
\bibitem{dok133} Pascher, J.: \emph{Fraktale Korrektur $K_\text{frak}$ (Dok.\,133)},
  T0-Theorie-Reihe, 2025.
\bibitem{dok165} Pascher, J.: \emph{Hubble-Verhältnis auf der $\xi$-Leiter (Dok.\,165)},
  T0-Theorie-Reihe, 2026.
\bibitem{dok166} Pascher, J.: \emph{Casimir-CMB-$H_0$-Brücke (Dok.\,166)},
  T0-Theorie-Reihe, 2026.
\bibitem{dok167} Pascher, J.: \emph{LiNbO$_3$ auf der $\xi$-Leiter (Dok.\,167)},
  T0-Theorie-Reihe, 2026.
\end{thebibliography}
